% !TEX root = ../main.tex
% !TEX program = lualatex
\documentclass[../main.tex]{subfiles}
\IfSubfilesClassLoaded{\externaldocument{\subfix{../build/main}}}{}
% ====================
% File: 13B_Contract_Administration.tex
% ====================
\begin{document}
\IfSubfilesClassLoaded{\chapter{Contract Administration}}{}

% === CONTRACT ADMINISTRATION (EDO 3.2.6) ===
\section{Contract Administration}

\subsection{Summary}
\begin{description}
  \item[\textbf{Role in the process.}] Contract administration is a post-award function that protects Government interests, monitors contractor performance, manages changes, resolves disputes, and closes out contracts~\autocite{edo-3-2-6-contract-administration-2025,FAR-42-302-Contract-Administration,FAR-43-103-Contract-Modifications,FAR-49-402-3-Default-Procedure}.
  \item[\textbf{Key organizations.}] \ac{dcma}, \ac{dcaa}, and \ac{dfas} provide contract management, audit, and payment support to keep contracts compliant and funded~\autocite{edo-3-2-6-contract-administration-2025,DCMA-About-2026,DCAA-About-2026,DFAS-Doing-Business-2026}.
  \item[\textbf{Key personnel.}] The \ac{pco}, \ac{aco}, \ac{cor}, and \ac{tco} coordinate with the \ac{pmo} and program integrator to manage contract execution~\autocite{edo-3-2-6-contract-administration-2025,FAR-42-302-Contract-Administration}.
  \item[\textbf{Changes and commitments.}] Distinguish change orders from supplemental agreements, track NECP/ECP actions, and avoid constructive changes and unauthorized commitments~\autocite{edo-3-2-6-contract-administration-2025}.
  \item[\textbf{Disputes and termination.}] Use \ac{adr} and the disputes process, and apply termination for convenience, cause, or default based on commercial/noncommercial context~\autocite{edo-3-2-6-contract-administration-2025,FAR-33-204-ADR-Policy,FAR-33-202-Disputes,FAR-49-402-3-Default-Procedure,FAR-52-212-4-Commercial-Terms}.
\end{description}

\subsection{Role in the Contracting Process}
\begin{itemize}
  \item Manage Government interests and monitor contractor quality, cost, and performance.
  \item Ensure contractor payment, administer contract modifications, and verify compliance with terms.
  \item Resolve disputes and execute contract close-out activities~\autocite{edo-3-2-6-contract-administration-2025}.
\end{itemize}

\subsection{Administration Organizations}
\begin{description}
  \item[\textbf{DCMA.}] Provides worldwide contract administration for non-shipbuilding contracts; functions include contract management, quality assurance, engineering/production oversight, property management, and program support~\autocite{edo-3-2-6-contract-administration-2025,DCMA-About-2026,FAR-42-302-Contract-Administration}.
  \item[\textbf{DCAA.}] Performs contract audit services and financial advisory support before and after award to inform negotiation, administration, and close-out actions~\autocite{edo-3-2-6-contract-administration-2025,DCAA-About-2026}.
  \item[\textbf{DFAS.}] Executes contract payments and must meet Prompt Payment Act timelines to protect contractor cash flow, especially for small businesses~\autocite{edo-3-2-6-contract-administration-2025,DFAS-Doing-Business-2026,FAR-32-904-Payment-Due-Dates}.
\end{description}

\begin{longtblr}
  [
    caption = {Contract administration organizations and services.},
    label   = {tab:contract-admin-orgs},
    entry   = {Contract Administration Organizations},
    remark{\tSource} = {\tabCite{edo-3-2-6-contract-administration-2025,DCMA-About-2026,DCAA-About-2026,DFAS-Doing-Business-2026,FAR-42-302-Contract-Administration}}
  ]{colspec = {@{} llll @{} }}
  \toprule
  {Organization} & {Program contract support}        & {Financial support}    & {Relationship to program} \\
  \midrule
  DCMA           & Full spectrum of contract support & -- & Assistance                \\
  DCAA           & Contract audit                    & -- & Audit                     \\
  DFAS           & Contract payments                 & Transaction processing & Functional                \\
  \bottomrule
\end{longtblr}

\note{The coursebook lifecycle diagram shows \ac{dcma} support from preaward planning through completion, including capability surveys, pricing support, production surveillance, acceptance, and closeout. Contract administration is the post-award phase; the agency's total support is broader~\autocite{edo-3-2-6-contract-administration-2025,DCMA-About-2026}.}

\subsection{Key Personnel and Relationships}
\begin{description}
  \item[\textbf{ACO and CAO.}] \ac{aco}s perform contract administration functions assigned by the \ac{pco} and operate within the \ac{cao}; \ac{supship} and \ac{rmc} equivalents perform \ac{aco} roles for shipbuilding and overhaul work~\autocite{edo-3-2-6-contract-administration-2025,FAR-42-302-Contract-Administration}.
  \item[\textbf{Program integrator.}] The program integrator serves as the \ac{cao} focal point for the program, leads the program support team, and provides status updates to the \ac{pmo} and \ac{dcma} chain of command~\autocite{edo-3-2-6-contract-administration-2025}.
  \item[\textbf{COR.}] The \ac{ko} appoints qualified personnel as \ac{cor}s with written authority; \ac{cor}s monitor and administer within that delegation, cannot change price, quality, quantity, delivery, or other contract terms, and unauthorized acts can create personal liability~\autocite{FAR-1-602-2-Responsibilities,edo-3-2-6-contract-administration-2025}.
  \item[\textbf{TCO.}] After a termination notice, the \ac{tco} negotiates an equitable settlement with the contractor~\autocite{edo-3-2-6-contract-administration-2025,FAR-42-302-Contract-Administration}.
\end{description}

\subsection{PMO and CAO Coordination}
\begin{itemize}
  \item Involve the \ac{cao} early, establish a \ac{moa}, and conduct a combined post-award conference.
  \item Inform the \ac{cao} of program developments, reviews (e.g., \ac{dab}), and funding issues; notify the \ac{cao} of planned visits~\autocite{FAR-42-402-CAO-Visits,edo-3-2-6-contract-administration-2025}.
\end{itemize}

\subsection{Cash Flow}
Timely Government payments drive contractor cash flow needed for materials, payroll, and subcontractor bills; late payments increase financing costs and can degrade performance, especially for small businesses. FAR~32.904 is the board-safe due-date anchor, while DFAS payment support includes invoice-status, PIEE, and MOCAS tools that help contractors and payment offices resolve payment workflow issues~\autocite{edo-3-2-6-contract-administration-2025,FAR-32-904-Payment-Due-Dates,DFAS-Doing-Business-2026}.

\subsection{Contract Changes and Modifications}
\begin{description}
  \item[\textbf{Change proposals.}] An \ac{necp} changes matters such as cost, delivery, or scope without changing configuration; an \ac{ecp} changes configuration or detailed technical specifications. Value engineering changes seek savings that may be shared under the applicable clause~\autocite{edo-3-2-6-contract-administration-2025}.
  \item[\textbf{Deviations and waivers.}] Deviations authorize pre-manufacture nonconformance; waivers accept nonconformance during or after manufacture as suitable-as-is, often with reduced cost~\autocite{edo-3-2-6-contract-administration-2025}.
  \item[\textbf{Modification requests.}] A field-origin \ac{fmrmod} is prepared in the \ac{cao} and approved by its designated authority; an \ac{hmr} originates with the program manager and receives the appropriate headquarters approval. Supporting \ac{necp}/\ac{ecp} packages and drawings support the warranted \ac{ko}'s modification action within delegated authority; program approval alone does not modify the contract~\autocite{edo-3-2-6-contract-administration-2025}.
  \item[\textbf{Contract changes.}] Supplemental agreements are bilateral modifications with equitable adjustments; change orders are unilateral and increase Government risk (undefinitized liabilities) while performance continues~\autocite{edo-3-2-6-contract-administration-2025,FAR-43-103-Contract-Modifications}.
\end{description}

\subsection{NECP/ECP Approval Levels}
The following is the coursebook shipbuilding hierarchy. Apply the current program's approval delegations and contract authority to an actual change~\autocite{edo-3-2-6-contract-administration-2025}.
\begin{description}
  \item[\textbf{Level 1 (OPNAV).}] Changes to military characteristics, technical specifications exceeding the \ac{cno} ship cost adjustment, schedule slips beyond contract delivery.
  \item[\textbf{Level 2 (SYSCOM).}] A reclama by an \ac{spm} \ac{ccb} member or change sponsor to the \ac{spm}'s change action; also changes affecting multiple program managers where a common decision cannot be reached.
  \item[\textbf{Level 3 (SPM).}] Field-activity or contractor changes affecting dates, guarantees, incentives, hull strength or safety, interfaces, or significant cost impacts.
  \item[\textbf{Level 4 (Field Activity).}] Changes not requiring higher approval~\autocite{edo-3-2-6-contract-administration-2025}.
\end{description}

\paragraph{SUPSHIP route.} In the coursebook flow, higher-level approval for Levels~1--3 proceeds through the \ac{spm}/\ac{hmr} route to the \ac{pco}/\ac{aco}; Level~4 uses the local board/field authority and \ac{fmrmod}. A warranted \ac{ko} then executes the appropriate change order or supplemental agreement. Approval of engineering, funding, or a modification request is distinct from execution of a contract modification~\autocite{edo-3-2-6-contract-administration-2025}.

\paragraph{RMC route.} A contractor's conditions-found report starts the coursebook repair workflow. The project team decides whether action is needed and whether the work is within original scope or constitutes growth/new work. Growth/new work requires the appropriate type-command approval and funding; an independent Government estimate and request for contract change support negotiation of the contractor's proposal. Resolve differences until price is fair and reasonable, document the agreement, and have the \ac{rmc} \ac{aco} execute the supplemental agreement. Reporting a condition or approving a technical solution does not itself authorize extra work~\autocite{edo-3-2-6-contract-administration-2025}.

\subsection{Constructive Changes and Unauthorized Commitments}
\begin{description}
  \item[\textbf{Constructive change.}] Government direction, conduct, or failure to act that a contractor reasonably treats as changing contract requirements; handle through the \ac{ko}, formal modification if appropriate, and equitable-adjustment or claim analysis rather than calling it ratification~\autocite{edo-3-2-6-contract-administration-2025,FAR-1-602-2-Responsibilities,FAR-43-103-Contract-Modifications}.
  \item[\textbf{Unauthorized commitment.}] An agreement made by a Government representative without authority; only a \ac{ko} acting within delegated authority can bind the Government, and unauthorized commitments are not binding unless ratified by an authorized official under \ac{far}~1.602-3~\autocite{FAR-1-602-3-Ratification,edo-3-2-6-contract-administration-2025}.
\end{description}

\subsection{Privity of Contract}
The Government has a direct contractual relationship with the prime contractor, not subcontractors; administration actions must respect this privity~\autocite{edo-3-2-6-contract-administration-2025}.

\subsection{Disputes and Alternative Dispute Resolution}
\begin{description}
  \item[\textbf{ADR.}] \ac{adr} promotes early, principled settlement of disputes and is encouraged to the maximum extent practicable before or alongside formal claim processes when appropriate~\autocite{FAR-33-204-ADR-Policy,edo-3-2-6-contract-administration-2025}.
  \item[\textbf{Disputes process.}] If unresolved at the contracting officer level, contractors may use the Contract Disputes Act process. After receipt of the \ac{ko}'s decision, the statutory alternatives generally are an agency-board appeal within 90 days or an action in the \ac{cofc} within 12 months; these are alternative routes, not mandatory sequential steps. Obtain counsel on the applicable filing and review rules~\autocite{USC-41-7104-Contract-Appeals,FAR-33-202-Disputes,edo-3-2-6-contract-administration-2025}.
\end{description}

\subsection{Termination Types and Notices}
\begin{description}
  \item[\textbf{Convenience.}] Used when requirements change, funds are insufficient, quantities shrink, or Government needs evolve; it is a unique Government right~\autocite{edo-3-2-6-contract-administration-2025}.
  \item[\textbf{Cause vs.\ default.}] Termination for cause applies to commercial contracts; termination for default applies to noncommercial contracts when performance fails~\autocite{edo-3-2-6-contract-administration-2025,FAR-52-212-4-Commercial-Terms}.
  \item[\textbf{Notices.}] For applicable progress or other performance failures, a cure notice allows 10 days or longer if necessary; failure to deliver on time may permit default without an advance cure notice. A show-cause notice seeks the contractor's explanation when appropriate. A termination notice states the action and consequences; it does not automatically grant a new cure period~\autocite{edo-3-2-6-contract-administration-2025,FAR-49-402-3-Default-Procedure,FAR-Overhaul-Part49-2026,FAR-49-607-Delinquency-Notices}.
\end{description}

\subsection{Commercial Item Definitions}
\begin{description}
  \item[\textbf{Commercial products and services.}] FAR 2.101 definitions cover products, services, and non-developmental items available in the commercial marketplace, including customary or minor modifications~\autocite{FAR-2-101-Definitions,edo-3-2-6-contract-administration-2025}.
  \item[\textbf{COTS.}] A \ac{cots} item is a qualifying commercial product sold in substantial quantities in the commercial marketplace and offered to the Government without modification in the same form; bulk cargo is excluded. The term does not classify all commercial services~\autocite{FAR-2-101-Definitions,edo-3-2-6-contract-administration-2025}.
\end{description}

\subsection{Government Obligation to Pay}
For cost-reimbursement contracts, recovery of incurred costs, necessary termination work, approved subcontract settlements, and any adjusted fee depends on the applicable termination clause and settlement rules. Do not assume entitlement to anticipated profit on unperformed work~\autocite{edo-3-2-6-contract-administration-2025,FAR-Overhaul-Part49-2026}.

An epidemic, strike, flood, Government act, embargo, fire, or unusually severe weather is an example of a potentially excusable cause, not automatic relief. Under the applicable clause, the failure must be beyond the contractor's control and without its fault or negligence; subcontractor failures have additional conditions. A justified schedule extension is distinct from entitlement to additional money~\autocite{FAR-52-249-14-Excusable-Delays,edo-3-2-6-contract-administration-2025}.

\subsection{Quick Review}
\begin{itemize}
  \item Contract administration starts after award and protects Government interests through performance surveillance, payment, modifications, dispute handling, termination, and closeout.
  \item \ac{dcma} manages administration support, \ac{dcaa} audits, and \ac{dfas} pays; for ship work, \ac{supship} and \ac{rmc} equivalents often perform administration roles.
  \item \acp{cor} monitor within written delegation only; they cannot change price, quantity, quality, delivery, or any other contract term.
  \item A supplemental agreement is bilateral; a change order is unilateral and creates undefinitized risk until the equitable adjustment is settled. FAR~43.103 is the clean official split for bilateral versus unilateral modifications.
  \item Constructive changes come from Government conduct that effectively changes the work; unauthorized commitments come from officials acting without authority and require ratification to bind the Government.
  \item Use \ac{adr} and the disputes process deliberately, and distinguish termination for convenience from commercial termination for cause and noncommercial termination for default based on contract context.
\end{itemize}

%====================
% End of file
%====================
\IfSubfilesClassLoaded{
  \RenewDocumentCommand{\entryname}{}{\textbf{\color{Modern} Acronym}}
  \RenewDocumentCommand{\descriptionname}{}{\textbf{\color{Modern} Definition}}
  \printnoidxglossary[
    type=\acronymtype,
    title=Acronyms,
    style=long-booktabs
  ]}{}
\end{document}
